\section{Die Trennung, auf die alles ankommt}

\textbf{[K]} Die Theorie sagt Verhältnisse voraus, nicht Absolutwerte.

Dieser Satz ist die wichtigste methodische Aussage des ganzen Blocks, und er
schneidet die Fehlerdiskussion in zwei Hälften, die nichts miteinander zu tun
haben.

Ein dimensionsloses Verhältnis zweier Größen derselben Ebene ist intrinsisch:
es hängt an der Geometrie und an nichts sonst. Es braucht keine Einheit, keine
Brückenkonstante, keinen Korrekturfaktor -- der kürzt sich heraus (A040).

Ein Absolutwert dagegen -- eine Masse in Megaelektronenvolt, eine Länge in
Metern -- verlangt zusätzlich eine Umrechnung, und diese Umrechnung ist nicht
Teil der Geometrie.

\textbf{[B]} Daraus folgt eine scharfe Diagnoseregel: \emph{jede Abweichung
zwischen Theorie und Messung ist zunächst daraufhin zu prüfen, ob sie im
geometrischen Kern oder in der Übersetzungsschicht sitzt.} Beides gleich zu
behandeln, verwischt die Aussage.

\section{Was die Naturkonstanten hier sind}

\textbf{[SETZUNG]} Im natürlichen System der Theorie sind $\hbar$, $c$ und
$k_B$ gleich Eins. Sie sind keine physikalischen Größen, sondern
Einheitenbrücken.

Das Wort Buchhaltung, das im Korpus dafür steht, wird regelmäßig als abwertend
missverstanden. Es ist das Gegenteil gemeint. Buchhaltung heißt hier
\emph{notwendig}, nicht \emph{entbehrlich}:

\begin{itemize}
  \item $c$ ist die Brücke zwischen Raum- und Zeiteinheit,
  \item $\hbar$ die Brücke zwischen Masse- und Zeiteinheit,
  \item $k_B$ die Brücke zwischen Temperatur und Energie.
\end{itemize}

\textbf{[B]} Das SI trennt Meter, Sekunde, Kilogramm und Kelvin in eigene
Einheiten. Sobald in diesem System \emph{gerechnet} wird, sind die
Umrechnungsfaktoren unverzichtbar -- ohne $k_B$ käme eine Energie in Kelvin
heraus. Ontologisch sind sie sekundär, operativ unentbehrlich.

\section{Die vier Kleidungen}

\textbf{[K]} Eine einzige Größe erscheint in vier Gestalten:
\[
  m_0 c^2 \;=\; E_0 \;=\; k_B T_0 \;=\; \hbar/\tilde T .
\]

Teilt man durch $c^2$, erhält man eine Masse; durch $\hbar$, eine Zeit; durch
$k_B$, eine Temperatur. Es sind vier Ausdrucksweisen desselben Sachverhalts,
nicht vier Größen. Diese Identität ist \emph{exakt} -- sie ist die Definition
der Umrechnung, und in ihr steckt keine Näherung.

\textbf{[B]} $E_0$ selbst ist eine empirische Brücke: $E_0=\sqrt{m_em_\mu}
\approx7{,}398\,\text{MeV}$ folgt aus gemessenen Leptonmassen, nicht aus $\xi$
allein. Die geometrische Herleitung steht in A130; die Einheitenprüfung in A085.
Wer nur dieses Dokument liest, muss wissen: $E_0$ ist hier als gegebene
Brückenkonstante geführt, nicht als abgeleitete Größe.

\textbf{[B]} Wo also entstehen Abweichungen? Nicht in dieser Identität. Sie
entstehen dort, wo eine theoretisch bestimmte dimensionslose Zahl gegen eine
gemessene Größe gehalten wird -- und dort sitzen sie in der Verkleidung, nicht
in der Umrechnung.

\section{Brückenkonstanten}

\textbf{[SETZUNG]} Um absolute Werte anzugeben, braucht die Theorie mindestens
eine Größe mit Dimension. Diese Brückenkonstante ist eine deklarierte Eingabe.

Das ist keine Schwäche des Ansatzes, sondern eine Notwendigkeit jeder Theorie,
die aus dimensionslosen Strukturen absolute Zahlen erzeugen will: aus reinen
Zahlen folgt keine Masse in Kilogramm. Entscheidend ist allein, dass die Zahl
der Brückenkonstanten offen genannt wird und klein bleibt.

\textbf{[B]} Zugleich gilt: eine Brückenkonstante zu wählen ist eine
Kalibrierung, kein Test. Eine Übereinstimmung, die durch die Wahl der Brücke
erzeugt wurde, darf nicht als Bestätigung gezählt werden. Diese Regel wird in
A130 an einem konkreten Fall angewandt, wo sie den Unterschied macht zwischen
einem Erfolg und einer Selbstbestätigung.

\section{Der Toleranzmaßstab in Einheiten}

\textbf{[B]} Aus der Trennung folgt, wie eine Abweichung zu lesen ist. Drei
Angaben gehören zusammen: der Wert, die Messgenauigkeit der Vergleichsgröße und
die Bestimmungstoleranz der Theoriegröße.

Eine Abweichung von einem Prozent kann bedeutungslos sein, wenn die
Theoriegröße nur auf zwei Stellen bestimmt ist -- und sie kann vernichtend sein,
wenn beide Seiten auf acht Stellen feststehen. Die Prozentzahl allein sagt
nichts. Diese Regel wird ab A100 durchgehend angewandt.

\section{Prüfskripte}

\begin{itemize}
  \item Exaktheit der Vier-Kleidungen-Identität $m_0c^2=E_0=k_BT_0=\hbar/\tilde T$
    in SI-Zahlen, mit Ausweis der relativen Abweichung (erwartet: Null bis auf
    Rundung): \texttt{python/A\_Serie\_Skripte/a080\_vier\_kleidungen.py}
  \item Invarianz von Verhältnissen unter Wechsel der Brückenkonstante, gegen
    die Nicht-Invarianz der Absolutwerte:
    \texttt{python/A\_Serie\_Skripte/a080\_verhaeltnis\_invarianz.py}
\end{itemize}

\section{Was offen bleibt}

\textbf{Erstens ist die Zahl der Brückenkonstanten hier nicht bilanziert.} Sie
gehört in die Gesamtbilanz A230, nicht hierher -- aber sie muss dort stehen,
sonst ist die Ein-Parameter-Aussage der Theorie nicht prüfbar.

\textbf{Zweitens ist nicht gezeigt, dass die Brückenkonstante aus der Geometrie
folgen könnte.} Solange sie deklariert wird, bleibt die absolute Skala eine
Eingabe. Das ist konsistent mit A020, wo dasselbe für die Größenordnung von
$\xi$ gilt, und es ist dieselbe Lücke in anderer Gestalt.

\section{Ersetzt}

Dieses Dokument nimmt auf: Dok.~013 (SI), Dok.~014 (natürlich/SI), Dok.~015
(Systematik der Natureinheiten), Dok.~122 (Verhältnis und Absolutes), Dok.~134
(Einheitenkonventionen), Dok.~261 (statische Natureinheiten). Eingearbeitet: P40 (Verhältnisse exakt, Abweichungen bei der Umrechnung).

\section{Verwendung von KI-Werkzeugen}

Dieses Dokument ist unter Verwendung KI-gestützter Werkzeuge entstanden. Die
Fragestellungen, die Setzungen und alle inhaltlichen Entscheidungen stammen vom
Autor; die Werkzeuge formulieren aus, rechnen nach und erstellen Prüfskripte.
Die Verantwortung für Inhalt und Fehler liegt vollständig beim Autor. Der
ausführliche Wortlaut steht in A010.
